% ======================================================================
% VOL III REARCHITECTURE PROPOSAL
% Post-CY-A_3 proof, post-K3 Yangian theorem, post-180-agent session
% April 2026
% ======================================================================

\documentclass[11pt]{article}
\usepackage[margin=1.2in]{geometry}
\usepackage{enumitem,booktabs,longtable}
\usepackage{amsmath,amssymb}
\usepackage[T1]{fontenc}
\frenchspacing

\newcommand{\CY}{\mathrm{CY}}
\newcommand{\Eone}{E_1}
\newcommand{\Etwo}{E_2}
\newcommand{\Einf}{E_\infty}
\newcommand{\En}{E_n}
\newcommand{\CoHA}{\mathrm{CoHA}}
\newcommand{\Rep}{\mathrm{Rep}}
\newcommand{\HH}{\mathrm{HH}}
\newcommand{\Ainf}{A_\infty}
\newcommand{\Phi}{\Phi}
\newcommand{\C}{\mathbb{C}}
\newcommand{\Z}{\mathbb{Z}}

\title{Volume III Rearchitecture Proposal:\\
\textit{Calabi--Yau Quantum Groups} after the proof of CY-A$_3$}
\author{Raeez Lorgat}
\date{April 2026}

\begin{document}
\maketitle

\tableofcontents
\newpage

% ======================================================================
\section{Diagnosis: what changed and why the old structure fails}
% ======================================================================

The current five-part structure was designed around the assumption that
CY-A$_3$ (the CY-to-chiral functor at $d=3$) was an open programme.
Four developments have made that structure obsolete:

\begin{enumerate}[label=(\arabic*)]
\item \textbf{CY-A$_3$ has a proved object-level core.} The
 $\infty$-categorical proof
 ($\HH^{-2}_{E_1}=0$, $E_3$-lifting contractible) resolves the local
 obstruction on verified framed loci. The output
 $\Phi_3^{(\Sigma_2,C)}(\mathcal C)$ exists there as an
 existence-and-rigidity theorem on objects. Arbitrary CY$_3$ morphism
 functoriality and the global quantum vertex group $G(\mathcal C)$
 remain outside the theorem. Part~I's narrative arc---``we prove
 $d=2$, we hope for $d=3$''---must be replaced by the sharper
 trichotomy: $d=2$ unconditional, $d=3$ verified object-level loci,
 morphism/global-$G$ conjectural.

\item \textbf{The K3 abelian Yangian is a theorem.} The complete
 6-part generators-and-relations presentation of $Y(\mathfrak{g}_{K3})$
 (24 Heisenberg generators $J_i$, Mukai-signature Serre relations,
 degree-$(24,24)$ structure function) is proved. This is too large for
 a section inside ``Toroidal and elliptic algebras''; it deserves a
 dedicated chapter.

\item \textbf{The $[m_3, B^{(2)}]$ saga.} The proof of the CY-A$_3$
 vanishing went through 3 wrong proofs and 4 independent computations
 before the TCFT/$\infty$-categorical resolution. This story is
 mathematically significant (it illuminates the \emph{method} of
 chain-level CY algebra) and pedagogically valuable. The current
 structure has no place for it.

\item \textbf{Six routes to $G(K3\times E)$.} The BKM, lattice VOA,
 Kummer, CY-to-chiral, sigma model, and BLLPR routes are now
 established as distinct constructions with proved pieces and
 conjectural comparison arrows. The current structure scatters these across
 \texttt{toroidal\_elliptic.tex} (7190 lines),
 \texttt{k3\_times\_e.tex} (5986 lines, commented out), and
 \texttt{modular\_koszul\_bridge.tex}. A unifying chapter is needed.
\end{enumerate}

Additional structural pressures:

\begin{itemize}
\item The CFG25 comparison (24\% of 3d CS structures lift to 6d)
 reframes Part~IV: the ``seven faces'' are no longer abstract---they
 are constrained by the 3d$\to$6d obstruction catalogue.
\item Part~V (``The CY Frontier'') currently contains only
 \texttt{geometric\_langlands.tex}. The frontier is no longer CY-A$_3$;
 it is the nonabelian Yangian, quantum toroidal algebras,
 root-of-unity phenomena, the ZTE correction, and the super-Yangian
 $Y(\mathfrak{gl}(4|20))$.
\item The $\kappa_{\mathrm{ch}}$ formula needs fixing (gives 24 for K3,
 should be 2). This affects the modular bridge chapter.
\item The shadow--Feynman dictionary ($L$-loop $= S_{L+1}$) and the
 class~M Borel summability result ($E_3$ bar $= 6^g$) have no natural
 home in the current structure.
\end{itemize}

% ======================================================================
\section{The new structure: seven parts}
% ======================================================================

The ideal Vol~III has \textbf{seven parts} and approximately
\textbf{28 chapters} (plus introduction, appendices, and bibliography).
The seven parts trace the logical arc:

\begin{center}
\renewcommand{\arraystretch}{1.15}
\begin{tabular}{cl}
\toprule
Part & Title \\
\midrule
I & Foundations: CY Categories and Cyclic $\Ainf$ \\
II & The CY-to-Chiral Functor \\
III & The $\En$ Hierarchy and Chiral Quantum Groups \\
IV & The K3 Yangian \\
V & The CY Landscape \\
VI & The Seven Faces of $r_{\CY}(z)$ \\
VII & Frontiers \\
\bottomrule
\end{tabular}
\end{center}

The critical innovation: Parts~I--II are the \emph{input}
(CY categories $\to$ chiral algebras), Part~III is the
\emph{algebraic machine} ($\En$ operads, quantum groups,
Drinfeld centers), Part~IV is the \emph{central theorem}
(the K3 Yangian as the prototype worked example, now with a
verified object-level CY-A$_3$ input), Part~V collects further examples, Part~VI
connects to the bar-cobar machine and physics, and Part~VII
addresses open problems.

% ======================================================================
\section{Part-by-part chapter listing}
% ======================================================================

% --------------------------------------------------
\subsection{Part I: Foundations --- CY Categories and Cyclic $\Ainf$}
\label{sec:part-i}
% --------------------------------------------------

\textit{Goal: establish the categorical input. All material here is
$d$-independent and unconditional.}

\begin{enumerate}[label=\textbf{Ch.\,\arabic*.},leftmargin=3.5em]
\item \textbf{Introduction.}
 \textit{Source:} Current \texttt{introduction.tex} (578 lines),
 rewritten to reflect the verified object-level status of CY-A$_3$.
 The introduction should present the two-stage object-level
 correspondence
 $\Phi_d^{(\Sigma,C)}\colon \CY_d\text{-}\mathrm{Cat}\dashrightarrow
 \En\text{-}\mathrm{ChirAlg}$
 as the central construction, while keeping arbitrary-morphism
 functoriality conjectural. The four main theorems CY-A through CY-D
 are stated with their current status: CY-A proved at $d=2$ and on
 verified framed $d=3$ loci, CY-B programme, CY-C conjectural,
 CY-D programme.

\item \textbf{Calabi--Yau categories.}
 \textit{Source:} Current \texttt{cy\_categories.tex} (207 lines).
 Expand to $\sim$400 lines. CY$_d$ categories, Serre functors,
 smooth proper conditions, the CY pairing, examples ($D^b(\mathrm{Coh}(X))$,
 $\mathrm{Fuk}(X)$, $\mathrm{MF}(W)$). MF(W) CY dimension is $n-2$
 (AP-CY17).

\item \textbf{Cyclic $\Ainf$-algebras and the Connes $B$-operator.}
 \textit{Source:} Current \texttt{cyclic\_ainf.tex} (206 lines).
 Expand to $\sim$400 lines. Cyclic $\Ainf$-structure, negative cyclic
 homology $\mathrm{HC}^-_d$, the $S^d$-framing, the Connes $B$-operator.
 The CY trace lives in $\mathrm{HC}^-_d$, NOT just $\HH_d\to k$ (AP-CY2).

\item \textbf{Hochschild calculus for CY categories.}
 \textit{Source:} Current \texttt{hochschild\_calculus.tex} (392 lines).
 Gerstenhaber bracket, Lie conformal algebra $\mathfrak{L}_\cC$,
 factorization envelope. The bridge from categorical Hochschild to
 chiral Hochschild (output of $\Phi$).
\end{enumerate}

\textit{Migration notes:} Chapters 2--4 are essentially the current
Part~I minus the E$_n$ and chiral material, which moves to Parts~II--III.

% --------------------------------------------------
\subsection{Part II: The CY-to-Chiral Functor}
\label{sec:part-ii}
% --------------------------------------------------

\textit{Goal: construct the object-level correspondence $\Phi_d$ at
honest scope. This is the load-bearing part. CY-A$_2$ is proved
unconditionally; CY-A$_3$ has a proved object-level core on verified
framed loci. Morphism functoriality and global $G(X)$ remain
conjectural.}

\begin{enumerate}[label=\textbf{Ch.\,\arabic*.},leftmargin=3.5em,start=5]
\item \textbf{The CY-to-chiral functor: construction and proof.}
 \textit{Source:} Current \texttt{cy\_to\_chiral.tex} (3882 lines),
 substantially rewritten. The four steps are presented uniformly as
 object-level constructions. CY-A$_2$ proof (existing). CY-A$_3$ proof:
 the $\infty$-categorical argument on the verified framed locus
 ($\HH^{-2}_{E_1}=0$, $E_3$-lifting contractible). The
 $\dashrightarrow$ on $\Phi_3^{(\Sigma_2,C)}$ becomes a solid arrow only
 inside that locus.

 Key structural change: the current chapter must distinguish three
 statuses: unconditional CY-A$_2$, verified object-level CY-A$_3$, and
 conjectural functorial/global CY-A$_3$. Conditional language remains
 load-bearing for CY-B, CY-C, CY-D, and for global $G(X)$.

\item \textbf{The $[m_3, B^{(2)}]$ saga: chain-level CY algebra.}
 \textit{Source:} NEW CHAPTER. $\sim$600--800 lines.
 This chapter tells the mathematical story of the $S^3$-framing
 obstruction: three wrong proofs (direct Hochschild, spectral sequence,
 model category), four independent computations (symbolic, numerical,
 TCFT, $\infty$-categorical), and the resolution. The pedagogical
 value: it exposes \emph{why} chain-level CY arguments are subtle,
 how the na\"ive $\HH^{-2}_{E_1}\neq 0$ expectation fails, and why
 the $\infty$-categorical framework was necessary.

 Sections: (i) The obstruction: $[m_3, B^{(2)}]$ and the failure of
 the direct approach. (ii) Three wrong proofs: what went wrong in each.
 (iii) The TCFT framework: how topological conformal field theory
 resolves the sign ambiguity. (iv) The $\infty$-categorical proof:
 $\HH^{-2}_{E_1}=0$ and the contractibility of the $E_3$-lifting
 space. (v) Lessons for chain-level algebra.

\item \textbf{The modular characteristic and CY Euler numbers.}
 \textit{Source:} Material from current \texttt{modular\_trace.tex}
 (156 lines) and the $\kappa_{\mathrm{ch}}$ sections of
 \texttt{cy\_to\_chiral.tex} and \texttt{modular\_koszul\_bridge.tex}.
 Consolidate. $\kappa_{\mathrm{ch}}=\chi^{\CY}$ at $d=2$ (proved).
 The $\kappa_{\mathrm{ch}}$ formula issue (gives 24, should be 2 for
 K3) is addressed. The $\kappa_\bullet$-spectrum is introduced here
 systematically.
\end{enumerate}

\textit{Migration notes:} Chapter~5 absorbs the core of the current
\texttt{cy\_to\_chiral.tex}. Chapter~6 is entirely new. Chapter~7
consolidates scattered $\kappa$ material from three current files.

% --------------------------------------------------
\subsection{Part III: The $\En$ Hierarchy and Chiral Quantum Groups}
\label{sec:part-iii}
% --------------------------------------------------

\textit{Goal: the algebraic infrastructure. $E_1$-chiral algebras,
$E_2$-chiral algebras, the Drinfeld center passage, quantum groups,
braided factorization, the $E_3$ bar complex, the $E_1$-chiral
bialgebra. This is the operadic heart of the volume.}

\begin{enumerate}[label=\textbf{Ch.\,\arabic*.},leftmargin=3.5em,start=8]
\item \textbf{$E_1$-chiral algebras and the ordered bar complex.}
 \textit{Source:} Current \texttt{e1\_chiral\_algebras.tex} (2023 lines).
 The Swiss-cheese operad, factorization on $\C\times\mathbb{R}$,
 the ordered bar $B^{\mathrm{ord}}$, the $E_1$-chiral bialgebra
 axioms (Section~7, $\sim$400 lines of new math from this session).

\item \textbf{$E_2$-chiral algebras and braided monoidal structure.}
 \textit{Source:} Current \texttt{e2\_chiral\_algebras.tex} (865 lines).
 The Deligne conjecture, $E_2$ from cyclic $\Ainf$, the three bar
 complexes $B_{E_1}$, $B_{E_2}$, $B_{E_\infty}$.

\item \textbf{$\En$-factorization and the stabilization theorem.}
 \textit{Source:} Current \texttt{en\_factorization.tex} (2288 lines).
 Dunn additivity, the framing obstruction tower, $E_1$-stabilization
 for $d\geq 3$, Bott periodicity. The $E_3$ bar cohomology:
 $(1+t)^{3g}$ for classes L,C; fails for class M (AP-CY21).

\item \textbf{Quantum groups: foundations.}
 \textit{Source:} Current \texttt{quantum\_groups\_foundations.tex}
 (260 lines). Drinfeld--Jimbo $U_q(\mathfrak{g})$, universal
 $R$-matrix, FRT/RTT, Kazhdan--Lusztig at root of unity.
 Expand to $\sim$500 lines.

\item \textbf{The Drinfeld center and the $E_1\to E_2$ passage.}
 \textit{Source:} Current \texttt{drinfeld\_center.tex} (1655 lines).
 Universal bulk, $\cZ(\Rep^{E_1}(A))\simeq\Rep^{E_2}(Z^{\mathrm{der}}_{\mathrm{ch}}(A))$.
 The K3 center obstruction ($>92\%$ non-local). The derived center
 (higher Deligne) for $E_2\to E_3$ promotion.

\item \textbf{Braided factorization and quantum vertex chiral groups.}
 \textit{Source:} Current \texttt{braided\_factorization.tex}
 (1361 lines) merged with current \texttt{quantum\_chiral\_algebras.tex}
 (2385 lines). The categorical $R$-matrix, the CY $R$-matrix
 construction, the target specification for $G(X)$. CY-A$_3$ supplies
 object-level $\Eone$-chiral shadows on verified loci; it does not
 construct $G(X)$ in general. The global quantum vertex chiral group is
 CY-C-level data, obstructed by the non-commutation of Drinfeld doubles
 with homotopy colimits.
\end{enumerate}

\textit{Migration notes:} This part absorbs the current Part~II's
quantum/braided material plus the $E_n$ chapter from Part~I. The
merge of \texttt{braided\_factorization.tex} and
\texttt{quantum\_chiral\_algebras.tex} eliminates redundancy between
these two chapters, which currently overlap in their treatment of
categorical $R$-matrices.

% --------------------------------------------------
\subsection{Part IV: The K3 Yangian}
\label{sec:part-iv}
% --------------------------------------------------

\textit{Goal: the central worked example, now a theorem. This part is
the mathematical climax of the volume. It presents the complete
construction $\mathrm{K3}\to Y(\mathfrak{g}_{K3})$ as the prototype
of the CY-to-quantum-group pipeline.}

\begin{enumerate}[label=\textbf{Ch.\,\arabic*.},leftmargin=3.5em,start=14]
\item \textbf{K3 surfaces and the Mukai lattice.}
 \textit{Source:} Opening sections of current \texttt{k3\_times\_e.tex}
 and \texttt{toroidal\_elliptic.tex}. The Mukai lattice
 $\Lambda^{4,20}$, the Mukai pairing, $\chi(\cO_{K3})=2$, the K3
 elliptic genus $\phi_{0,1}$.

\item \textbf{$\Phi_2(K3)=H_{\mathrm{Muk}}$: the CY$_2$ functor
 on K3.}
 \textit{Source:} NEW CHAPTER. $\sim$500--700 lines.
 The rigorous 4-step proof that $\Phi$ applied to $D^b(\mathrm{Coh}(K3))$
 produces the Mukai-lattice Heisenberg VOA. This is the first
 nontrivial complete computation of $\Phi$ for a compact manifold.

\item \textbf{The K3 double current algebra and the abelian Yangian.}
 \textit{Source:} Sections of current \texttt{k3\_times\_e.tex}
 (the double current algebra, Maulik--Okounkov $R$-matrix) plus
 the new 6-part generators-and-relations theorem. $\sim$800--1000 lines.
 The complete presentation of $Y(\mathfrak{g}_{K3})$: 24 Heisenberg
 generators, Mukai-signature Serre relations, degree-$(24,24)$
 structure function, the super-Yangian $Y(\mathfrak{gl}(4|20))$.

\item \textbf{The $K3\times E$ tower and the Igusa cusp form.}
 \textit{Source:} Bulk of current \texttt{k3\_times\_e.tex} (5986 lines)
 and \texttt{toroidal\_elliptic.tex} (the K3$\times$E-specific material,
 not the pure toroidal/elliptic algebra material). The DT zeta function,
 $\Lambda^{3,2}$, the BKM superalgebra $\mathfrak{g}_{\Delta_5}$,
 and the total-space four-construction spectrum
 $\{0,3,5,24\}$.

 The $\Delta_5$ denominator identity is theorem-grade on its
 Borcherds/Gritsenko input. Its identification with a bar Euler product
 is a conditional comparison statement requiring the object-level
 CY-A$_3$ input, the Vol~I bar-cobar identification, and the still
 conjectural CY-C comparison arrows.

\item \textbf{Six routes to $G(K3\times E)$.}
 \textit{Source:} NEW CHAPTER. $\sim$600--800 lines. A unifying
 treatment of the six independent constructions of the quantum vertex
 chiral group:
 \begin{enumerate}[label=(\roman*)]
  \item BKM route (Gritsenko--Nikulin denominator identity),
  \item Lattice VOA route ($V_\Lambda$ with $\Lambda$ the Mukai lattice),
  \item Kummer route (Steps 1--4 proved, Step 5 only on the verified
   object-level CY-A$_3$ locus),
  \item CY-to-chiral route
   ($\Phi_3^{(\Sigma_2,C)}$ applied objectwise to
   $D^b(\mathrm{Coh}(K3\times E))$),
  \item Sigma model route (chiral de Rham complex),
  \item BLLPR route (Bruinier--Laures--van der Geer--Zagier--Raum).
 \end{enumerate}
 The six routes are six \emph{different mathematical constructions}
 (only route~(iv) uses the CY-to-chiral correspondence; AP-CY60).
 Their conjectural convergence (Conjecture~CY-C) forms a commutative diagram,
 and the chapter traces which arrows are proved and which remain
 conjectural.

\item \textbf{The Kummer route: orbifold, excision, and resolution.}
 \textit{Source:} Material currently in \texttt{cy\_to\_chiral.tex}
 and \texttt{toroidal\_elliptic.tex} on the Kummer construction.
 $\sim$400--600 lines. Steps 1--4 are proved
 (\texttt{prop:kummer-orbifold}): $\int_{T^4} H_1 =$
 rank-16 Heisenberg, $\Z_2$-invariant $=$ rank-8, 16 twisted sectors,
 orbifold $\kappa_{\mathrm{ch}}=2$, resolution $8+32-16=24$.
 Step~5 follows only after the Kummer output is placed inside the
 verified object-level CY-A$_3$ locus; the global $G(K3\times E)$
 comparison remains CY-C-level data.
\end{enumerate}

\textit{Migration notes:} This part restructures the two massive
current example chapters (\texttt{toroidal\_elliptic.tex} at 7190 lines
and \texttt{k3\_times\_e.tex} at 5986 lines) into focused,
theorem-driven chapters. The pure toroidal/elliptic algebra material
(quantum toroidal $U_{q,t}$, Miki automorphism, double affine Hecke)
moves to Part~V.

% --------------------------------------------------
\subsection{Part V: The CY Landscape}
\label{sec:part-v}
% --------------------------------------------------

\textit{Goal: the remaining examples, organized by CY dimension and
type.}

\begin{enumerate}[label=\textbf{Ch.\,\arabic*.},leftmargin=3.5em,start=20]
\item \textbf{$\C^3$ and the affine Yangian $Y(\widehat{\mathfrak{gl}}_1)$.}
 \textit{Source:} Sections of current \texttt{toric\_cy3\_coha.tex}
 (1002 lines) specific to $\C^3$. The Jordan quiver, CoHA,
 $\cW_{1+\infty}$ identification, shadow tower (class M, infinite depth).

\item \textbf{Toric CY$_3$ and affine super Yangians.}
 \textit{Source:} Remainder of current \texttt{toric\_cy3\_coha.tex}.
 General toric CY$_3$ without compact 4-cycles (RSYZ), resolved
 conifold, local $\mathbb{P}^2$ (class M, not class L; AP-CY12).

\item \textbf{Toroidal and elliptic algebras.}
 \textit{Source:} The non-K3 material from current
 \texttt{toroidal\_elliptic.tex}. Quantum toroidal
 $U_{q,t}(\widehat{\widehat{\mathfrak{gl}}}_1)$, Miki automorphism
 (algebra-specific, not operadic; AP-CY22), elliptic Hall algebra,
 double affine Hecke algebras. $\sim$2000 lines.

\item \textbf{Fukaya categories and mirror symmetry.}
 \textit{Source:} Current \texttt{fukaya\_categories.tex} (710 lines).
 Lagrangian submanifolds, the $\Ainf$-structure on $\mathrm{Fuk}(X)$,
 HMS as input to $\Phi$.

\item \textbf{Derived categories, matrix factorizations, and
 non-compact CY.}
 \textit{Source:} Current \texttt{derived\_categories\_cy.tex}
 (290 lines) and \texttt{matrix\_factorizations.tex} (289 lines),
 merged. MF(W) CY dimension $= n-2$ (AP-CY17). The quintic:
 chart-gluing for non-toric CY$_3$.

\item \textbf{Quantum group representations from CY categories.}
 \textit{Source:} Current \texttt{quantum\_group\_reps.tex} (586 lines).
 Representation categories $\Rep^{E_2}$, modular tensor categories at
 root of unity. AP-CY5: Kazhdan--Lusztig requires root of unity.
\end{enumerate}

\textit{Migration notes:} The current Part~III (``CY Landscape'') is
restructured: the K3 material elevates to Part~IV, the toric/$\C^3$
material splits into two focused chapters, and the toroidal/elliptic
chapter is trimmed of its K3$\times$E content.

% --------------------------------------------------
\subsection{Part VI: The Seven Faces of $r_{\CY}(z)$}
\label{sec:part-vi}
% --------------------------------------------------

\textit{Goal: the connections to the bar-cobar machine (Vols~I--II),
to physics, and to other mathematical programmes.}

\begin{enumerate}[label=\textbf{Ch.\,\arabic*.},leftmargin=3.5em,start=26]
\item \textbf{The bar-cobar bridge to Volumes~I--II.}
 \textit{Source:} Current \texttt{bar\_cobar\_bridge.tex} (667 lines)
 expanded with the shadow--Feynman dictionary ($L$-loop $= S_{L+1}$)
 and the shadow--$\Ainf$-coproduct identification. The class~M Borel
 summability result ($E_3$ bar $= 6^g$). Cross-volume references
 to Vol~I \texttt{higher\_genus\_complementarity.tex} and Vol~II
 \texttt{foundations\_recast\_draft.tex}.

\item \textbf{The seven faces of $r_{\CY}(z)$.}
 \textit{Source:} Current \texttt{cy\_holographic\_datum\_master.tex}
 (1014 lines). Now reframed by the CFG25 comparison: only 24\% of 3d
 CS structures lift to 6d. The seven faces are constrained by the
 3d$\to$6d obstruction catalogue. New section: BKM Serre concentrated
 at $D=3$ (182 generators, finitely determined).

\item \textbf{Modular Koszul duality and CY geometry.}
 \textit{Source:} Current \texttt{modular\_koszul\_bridge.tex}
 (552 lines). CY modular convolution algebra, complementarity
 transport, CohFT on $\overline{\cM}_{g,n}$. K3 Koszul conductor $= 0$
 (free-field/KM branch).
\end{enumerate}

% --------------------------------------------------
\subsection{Part VII: Frontiers}
\label{sec:part-vii}
% --------------------------------------------------

\textit{Goal: the open problems and research directions. Now that the
object-level CY-A$_3$ obstruction is resolved on verified framed loci,
the frontier moves to functoriality, global $G(X)$, and the CY-B/C/D
comparisons.}

\begin{enumerate}[label=\textbf{Ch.\,\arabic*.},leftmargin=3.5em,start=29]
\item \textbf{The nonabelian Yangian and quantum toroidal algebras.}
 \textit{Source:} NEW CHAPTER. $\sim$400--600 lines. The passage
 from the abelian $Y(\mathfrak{g}_{K3})$ (proved) to the full
 nonabelian Yangian. Matrix Miura, Serre constraints from $\mathfrak{sl}_2$
 traces. Quantum toroidal from 3-chart atlas. Status: programme.

\item \textbf{The Zamolodchikov tetrahedron and $E_3$ corrections.}
 \textit{Source:} NEW CHAPTER. $\sim$400--600 lines. The computed
 negative result: Yang $R$-matrix does NOT satisfy ZTE. The $O(\kappa^2)$
 obstruction. ZTE correction EXISTS (deformation cohomology). The
 $E_3$ is genuinely nontrivial. This is new mathematics.

\item \textbf{Geometric Langlands and CY quantum groups.}
 \textit{Source:} Current \texttt{geometric\_langlands.tex}
 (629 lines). Feigin--Frenkel center, D-modules on $\mathrm{Bun}_G$,
 CY Langlands duality. Derived Satake from chiral data.

\item \textbf{Root of unity, categorification, and the CY frontier.}
 \textit{Source:} NEW CHAPTER. $\sim$400--600 lines. Root-of-unity
 phenomena for CY quantum groups, Khovanov-style categorification,
 tropical cluster structures, chiral Verlinde formula, the
 super-Yangian $Y(\mathfrak{gl}(4|20))$, and $\mathrm{Sp}_4(\Z)$
 modularity from the $E_3$ $S$-matrix.
\end{enumerate}


% ======================================================================
\section{Chapter-level migration table}
% ======================================================================

\begin{longtable}{p{5.5cm}p{1.3cm}p{5.5cm}p{1.3cm}}
\toprule
\textbf{Current file} & \textbf{Lines} & \textbf{New location} & \textbf{Part} \\
\midrule
\endhead
\texttt{introduction.tex} & 578 & Ch.\,1 (rewritten) & I \\
\texttt{cy\_categories.tex} & 207 & Ch.\,2 (expanded) & I \\
\texttt{cyclic\_ainf.tex} & 206 & Ch.\,3 (expanded) & I \\
\texttt{hochschild\_calculus.tex} & 392 & Ch.\,4 & I \\
\texttt{cy\_to\_chiral.tex} & 3882 & Ch.\,5 (rewritten) + Ch.\,7 + Ch.\,19 & II,IV \\
\texttt{modular\_trace.tex} & 156 & Ch.\,7 (merged) & II \\
\texttt{e1\_chiral\_algebras.tex} & 2023 & Ch.\,8 & III \\
\texttt{e2\_chiral\_algebras.tex} & 865 & Ch.\,9 & III \\
\texttt{en\_factorization.tex} & 2288 & Ch.\,10 & III \\
\texttt{quantum\_groups\_fdn.tex} & 260 & Ch.\,11 (expanded) & III \\
\texttt{drinfeld\_center.tex} & 1655 & Ch.\,12 & III \\
\texttt{quantum\_chiral\_alg.tex} & 2385 & Ch.\,13 (merged with braided) & III \\
\texttt{braided\_factorization.tex} & 1361 & Ch.\,13 (merged) & III \\
\texttt{toroidal\_elliptic.tex} & 7190 & Ch.\,14,17,22 (split) & IV,V \\
\texttt{k3\_times\_e.tex} & 5986 & Ch.\,14--19 (split) & IV \\
\texttt{toric\_cy3\_coha.tex} & 1002 & Ch.\,20--21 (split) & V \\
\texttt{fukaya\_categories.tex} & 710 & Ch.\,23 & V \\
\texttt{derived\_categories\_cy.tex} & 290 & Ch.\,24 (merged) & V \\
\texttt{matrix\_factorizations.tex} & 289 & Ch.\,24 (merged) & V \\
\texttt{quantum\_group\_reps.tex} & 586 & Ch.\,25 & V \\
\texttt{bar\_cobar\_bridge.tex} & 667 & Ch.\,26 (expanded) & VI \\
\texttt{cy\_holo\_datum\_master.tex} & 1014 & Ch.\,27 (reframed) & VI \\
\texttt{modular\_koszul\_bridge.tex} & 552 & Ch.\,28 & VI \\
\texttt{geometric\_langlands.tex} & 629 & Ch.\,31 & VII \\
--- & --- & Ch.\,6 ($[m_3,B^{(2)}]$ saga) & II \\
--- & --- & Ch.\,15 ($\Phi_2(K3)$) & IV \\
--- & --- & Ch.\,16 (K3 abelian Yangian) & IV \\
--- & --- & Ch.\,18 (six routes) & IV \\
--- & --- & Ch.\,29 (nonabelian Yangian) & VII \\
--- & --- & Ch.\,30 (ZTE correction) & VII \\
--- & --- & Ch.\,32 (root of unity/frontier) & VII \\
\bottomrule
\end{longtable}

\textit{New chapters:} 7 (Chs.\,6, 15, 16, 18, 29, 30, 32).\\
\textit{Merged chapters:} 3 merges (braided+quantum$\to$Ch.\,13;
derived+MF$\to$Ch.\,24; modular\_trace into Ch.\,7).\\
\textit{Split chapters:} 2 splits (toroidal\_elliptic $\to$ 3 chapters;
k3\_times\_e $\to$ absorbed into Part IV).

% ======================================================================
\section{The abstract}
% ======================================================================

\noindent\textit{Proposed abstract for the rearchitected volume:}

\medskip

\noindent
Volumes~I and~II construct the ordered bar complex
$B^{\mathrm{ord}}(\mathcal{A}) = T^c(s^{-1}\bar{\mathcal{A}})$
of any chiral algebra $\mathcal{A}$ and extract its modular invariants.
This volume constructs the geometric input: an object-level
correspondence
$\Phi_d^{(\Sigma,C)}\colon \CY_d\text{-}\mathrm{Cat}\dashrightarrow
\En\text{-}\mathrm{ChirAlg}$
from Calabi--Yau categories of dimension $d$ to $\En$-chiral algebras,
proved unconditionally at $d=2$ and on verified framed loci at $d=3$.

A $\CY_d$ category $\mathcal{C}$ carries a cyclic $\Ainf$-structure
whose $S^d$-framing determines an operadic enhancement. The
correspondence $\Phi_d^{(\Sigma,C)}$ converts this datum into a chiral algebra on a curve in four
steps: cyclic $\Ainf$-structure $\to$ Lie conformal algebra $\to$
factorization envelope $\to$ $S^d$-enhanced quantization. For $d=2$,
the output is an $\Etwo$-chiral algebra (Theorem CY-A$_2$); for
$d=3$, the chain-level $S^3$-framing obstruction is resolved by an
$\infty$-categorical argument on the verified framed locus
($\HH^{-2}_{E_1}=0$), producing an $\Eone$-chiral algebra there
(Theorem CY-A$_3$).

The central application is the K3 Yangian. The proved $d=2$
correspondence applied to the K3 surface produces the Mukai-lattice Heisenberg VOA
($\Phi_2(K3)=H_{\mathrm{Muk}}$, Theorem). The K3 double current
algebra $\mathfrak{g}_{K3}$ quantizes to an abelian Yangian
$Y(\mathfrak{g}_{K3})$ with 24 generators, Mukai-signature Serre
relations, and degree-$(24,24)$ structure function (Theorem). For the
threefold $K3\times E$, six independent routes---BKM, lattice VOA,
Kummer, CY-to-chiral, sigma model, BLLPR---define a candidate quantum
vertex chiral group $G(K3\times E)$. The BKM route is governed by the
\(\Delta_5\) denominator form, equivalently the square-root convention
for the Igusa form \(\Phi_{10}\); the comparison arrows among routes
remain CY-C-level data.

The $\kappa_\bullet$-spectrum
$\mathrm{Spec}_{\kappa_\bullet}(K3\times E) = \{0,3,5,24\}$ records
the total-space four-construction values; the value \(2\) is the
separate K3 fibre categorical Euler witness. The $\En$-hierarchy ($\Einf$ at $d=1$, $\Etwo$ at
$d=2$, $\Eone$ at $d\geq 3$) is controlled by an effective framing
obstruction $\mathrm{Obs}_{\mathrm{eff}}(d)\in\pi_d(BU)$, trivial
when $d\bmod 8\in\{1,3,7\}$. The $E_3$-bar cohomology is
$(1+t)^{3g}$ for shadow classes L and C, and $6^g$ (Borel summable)
for class M. The Zamolodchikov tetrahedron equation fails at
$O(\kappa^2)$ for the Yang $R$-matrix, proving that $E_3$ structure
is genuinely nontrivial beyond pairwise factorization.

% ======================================================================
\section{The preface: key paragraphs}
% ======================================================================

\noindent\textit{The following replaces the opening paragraphs of
the preface. The existing structural exposition (the $\En$ spiral,
the $\kappa_\bullet$-spectrum, the K3$\times$E prototype) is retained
but reframed around the verified object-level status of CY-A$_3$.}

\medskip

\noindent
A chiral algebra is a solution. What is the problem it solves?

Volumes~I and~II develop the bar-cobar machine that extracts modular
invariants from any chiral algebra~$A$ on a curve~$X$. Neither volume
constructs a single chiral algebra. This volume constructs the input:
both the chiral algebras and, from Calabi--Yau geometry, the concrete
$E_1$-chiral quantum groups that are the flesh on Volume~I's algebraic
skeleton.

The construction is an object-level correspondence
$\Phi_d^{(\Sigma,C)}\colon \CY_d\text{-}\mathrm{Cat}\dashrightarrow
\En\text{-}\mathrm{ChirAlg}$.
At $d=2$, the correspondence is proved as a theorem. At $d=3$, the
chain-level $S^3$-framing obstruction---the question of whether the
cyclic bar complex of a CY$_3$ category admits the correct higher
operadic enhancement---is resolved on the verified framed locus. The
$\infty$-categorical proof (Chapter~6) shows that the obstruction
group vanishes in that locus and that the relevant lifting space is
contractible. Global morphism functoriality remains conjectural.

With the object-level $\Phi_3$ in hand, the volume's centre of gravity shifts from
construction to application. The K3 Yangian---a 24-generator abelian
Yangian $Y(\mathfrak{g}_{K3})$ with Mukai-signature Serre relations
and degree-$(24,24)$ structure function---is the first nontrivial
quantum group constructed from Calabi--Yau geometry via $\Phi_2$.
For $K3\times E$, six independent constructions---of which only
Route~4 uses the CY-to-chiral correspondence (AP-CY60)---approach a
candidate common quantum vertex chiral group. The BKM denominator route is controlled by
\(\Delta_5\), with \(\Phi_{10}\) appearing through the square relation.
Their conjectural convergence (Conjecture~CY-C) is the strongest
evidence yet for the programme's coherence; the remaining open
problems (CY-B, CY-C, CY-D) concern the bar-cobar,
representation-theoretic, and modular properties of these newly
constructed objects.

\medskip

\noindent\textit{The remainder of the preface continues with the
$\kappa_\bullet$-spectrum, the $\En$-hierarchy spiral, and the K3$\times$E
prototype, as in the current version but with CY-A$_3$ passages split
between proved object-level loci and conjectural functorial/global
extensions.}


% ======================================================================
\section{The introduction: structural outline}
% ======================================================================

The introduction (Chapter 1) should be rewritten to serve the new
seven-part structure. Key changes:

\begin{enumerate}[label=(\arabic*)]
\item \textbf{Opening paragraph.} State $\Phi_d^{(\Sigma,C)}$ as an
 object-level theorem at $d=2$ and on verified framed loci at $d=3$.
 The dashed arrow $\dashrightarrow$ on $\Phi_3$ becomes a solid arrow
 only inside those loci.

\item \textbf{Four theorems.} Update the status table:
 \begin{center}
 \begin{tabular}{lll}
 \toprule
 Theorem & Old status & New status \\
 \midrule
 CY-A (CY-to-chiral correspondence) & $d=2$ proved, $d=3$ programme & \textbf{$d=2$ proved; $d=3$ object-level theorem on verified loci; morphisms conjectural} \\
 CY-B ($E_n$-chiral Koszul) & Programme & Programme \\
 CY-C (Quantum group realization) & Conjectural & Conjectural \\
 CY-D (Modular CY characteristic) & Programme & Programme \\
 \bottomrule
 \end{tabular}
 \end{center}

\item \textbf{Guide to the reader.} Provide three reading paths:
 \begin{itemize}
 \item \textit{The theorem path} (Parts I--II): for readers who want
  the construction of $\Phi$ and the proof of CY-A$_3$.
 \item \textit{The K3 path} (Part IV): for readers interested in the
  K3 Yangian and K3$\times$E as a self-contained story.
 \item \textit{The landscape path} (Part V$\to$VI): for readers coming
  from enumerative geometry, CoHA, or physics.
 \end{itemize}

\item \textbf{Connective tissue.} A diagram showing the logical
 dependencies between parts:
 \[
 \text{I} \to \text{II} \to \text{III} \to
 \begin{cases} \text{IV (K3)} \\ \text{V (Landscape)} \end{cases}
 \to \text{VI} \to \text{VII}
 \]
 Parts IV and V are independent of each other. Part VI draws on both.
 Part VII is self-contained but references all prior parts.
\end{enumerate}


% ======================================================================
\section{Connective tissue between chapters}
% ======================================================================

\subsection{Part I $\to$ Part II}

The bridge is the Lie conformal algebra $\mathfrak{L}_\cC$. Part~I
constructs $\HH^\bullet(\cC)$ with its Gerstenhaber bracket and CY
pairing; Part~II constructs the object-level $\Phi_d^{(\Sigma,C)}$
from these data. The handoff is Chapter~4 (Hochschild calculus)
$\to$ Chapter~5 (CY-to-chiral correspondence): the former builds the
Lie conformal algebra, the latter applies the factorization envelope.

\subsection{Part II $\to$ Part III}

The bridge is the $S^d$-framing. Chapter~5 proves the object-level
existence statement on the verified loci; Chapter~6 tells the story of
\emph{why} the proof was hard. Part~III then develops the operadic
infrastructure ($E_1$/$E_2$/Drinfeld center) that organizes the output
of $\Phi_d^{(\Sigma,C)}$. The logical dependency: Part~III does not
depend on the functorial CY-A$_3$ conjecture (the operadic theory is
independent), but the \emph{motivation} comes from Part~II.

\subsection{Part III $\to$ Part IV}

The bridge is the K3 Yangian as the first complete instantiation of
the $\Eone$-chiral bialgebra axioms. Chapter~13 (braided factorization)
sets up the target specification for $G(X)$; Chapter~14 (K3 surfaces)
provides the geometric input; Chapters~15--19 work through the
computation. The $E_1$-chiral bialgebra axioms of Chapter~8 are
\emph{verified} for the K3 Yangian in Chapter~16.

\subsection{Part IV $\to$ Part V}

The bridge is the passage from the K3 prototype to the general
landscape. Chapter~18 (six routes) shows that $K3\times E$ is
over-determined; Part~V shows that other geometries ($\C^3$, toric,
quintic) have fewer routes and more open questions. The contrast
motivates the frontier (Part~VII).

\subsection{Part V $\to$ Part VI}

The bridge is the bar-cobar complex. Part~V computes chiral algebras
for specific CY geometries; Part~VI feeds those algebras into the
Vol~I machine ($B^{\mathrm{ord}}$, $\Theta_A$, shadow tower,
Koszul duality). The shadow--Feynman dictionary ($L$-loop $= S_{L+1}$)
is the main new cross-volume result.

\subsection{Part VI $\to$ Part VII}

The bridge is the list of open problems generated by the
completed programme. The seven faces are now constrained (CFG25: 24\%
lift to 6d); the Zamolodchikov tetrahedron shows $E_3$ is nontrivial;
the nonabelian Yangian is the next natural target. Each frontier chapter
begins by stating which results from Parts~I--VI it relies on.


% ======================================================================
\section{Indicative sizes (no page-count target)}
% ======================================================================

The numbers below are working estimates of the relative weight each
Part bears in the seven-part rearchitecture. They are not page-count
targets. Volume~III's length is whatever it must be to work out the
inner music of the subject: the $\En$ hierarchy at every CY dimension,
the K3 Yangian in full, the CY landscape, and the frontier programmes.
A short Part is short because its content is concentrated; a long Part
is long because its content demands the room.

\begin{center}
\renewcommand{\arraystretch}{1.1}
\begin{tabular}{clr}
\toprule
Part & Title & Indicative lines \\
\midrule
I & Foundations & 1800 \\
II & CY-to-Chiral Functor & 5000 \\
III & $\En$ Hierarchy \& Chiral QG & 9000 \\
IV & The K3 Yangian & 9000 \\
V & CY Landscape & 5500 \\
VI & Seven Faces of $r_{\CY}(z)$ & 3000 \\
VII & Frontiers & 2500 \\
\midrule
 & Appendices, bibliography, front matter & 1500 \\
\bottomrule
\end{tabular}
\end{center}

The increase from the current $\sim$371pp baseline is concentrated in
Part~IV (the K3 Yangian, elevating from subsections to a full Part)
and Part~II (the $[m_3, B^{(2)}]$ saga). The realised length will
exceed every prior estimate as the inner mathematics is worked out;
that is the design.


% ======================================================================
\section{CLAUDE.md roadmap update}
% ======================================================================

The CLAUDE.md roadmap section should be updated to reflect the new
structure. The key changes:

\begin{enumerate}[label=(\arabic*)]
\item CY-A$_3$: status changes from programme to proved object-level on
 verified framed H1--H4 specialisation loci; arbitrary CY$_3$ morphism
 functoriality and global $G(\mathcal C)$ remain conjectural.
\item The five ``load-bearing open problems'' list is updated:
 CY-A$_3$ is removed; the nonabelian Yangian, ZTE correction,
 root-of-unity, CY-B (Koszul), and Sp$_4(\Z)$ modularity remain.
\item The ``Status by dimension'' table is updated:
 $d=3$ changes from programme to framed object-level theorem; morphism
 functoriality remains conjectural.
\item The compute engine list is updated with the new engines.
\item The architecture diagram is updated with the seven-part structure.
\end{enumerate}


% ======================================================================
\section{Implementation plan}
% ======================================================================

The rearchitecture can be implemented incrementally without breaking
the build at any step:

\begin{enumerate}[label=\textbf{Phase \arabic*.}]
\item \textbf{Reorganize main.tex} (1 session). Update the
 \texttt{\textbackslash{}part} and \texttt{\textbackslash{}input}
 structure. This is the skeleton change. All existing chapter files
 continue to compile; new chapters are initially stubs with
 \texttt{\textbackslash{}chapter\{...\}} and a one-paragraph summary.

\item \textbf{Write new chapters} (3--5 sessions). Chapters 6, 15, 16,
 18, 29, 30, 32. Each is 400--800 lines. Priority order: Ch.\,6
 ($[m_3,B^{(2)}]$), Ch.\,16 (K3 Yangian theorem), Ch.\,18 (six routes),
 Ch.\,15 ($\Phi_2(K3)$), Ch.\,30 (ZTE), Ch.\,29 (nonabelian Yangian),
 Ch.\,32 (frontier).

\item \textbf{Rewrite Chapter 5} (1--2 sessions). The CY-to-chiral
 chapter needs the CY-A$_3$ object-level proof integrated while
 preserving conditional language for morphism functoriality and global
 $G(X)$.

\item \textbf{Split the megachapters} (1--2 sessions).
 \texttt{toroidal\_elliptic.tex} (7190 lines) and
 \texttt{k3\_times\_e.tex} (5986 lines) are divided along the
 boundaries described above.

\item \textbf{Merge redundant chapters} (1 session).
 \texttt{braided\_factorization.tex} + \texttt{quantum\_chiral\_algebras.tex};
 \texttt{derived\_categories\_cy.tex} + \texttt{matrix\_factorizations.tex};
 \texttt{modular\_trace.tex} into Chapter~7.

\item \textbf{Rewrite abstract, preface, introduction} (1 session).
 Using the drafts in Sections~5--7 of this proposal.

\item \textbf{Cross-reference audit} (1 session). Update all
 \texttt{\textbackslash{}ref\{part:...\}} labels. Grep all three volumes
 for stale Part references (AP-CY13).

\item \textbf{Test suite verification} (continuous). The 19,838+ tests
 must pass after each phase. No test is removed; new tests are added
 for new chapters.
\end{enumerate}


\end{document}
